\documentclass[12pt]{article}
\usepackage{amsmath, amssymb}
\usepackage[margin=1in]{geometry}
\setlength{\parskip}{6pt}

\title{QUBO Formulation: Score-Based Bayesian Network Structure Learning with Candidate Parent Sets}
\date{}

\begin{document}
\maketitle

\subsection*{Score-Based Bayesian Network Structure Learning with Candidate Parent Sets}

\paragraph{Problem Statement.}
We are given a set of $N$ nodes indexed by $i \in \{0, 1, \dots, N-1\}$ with a fixed topological order $0 < 1 < \dots < N-1$. For each node $i$, there exists a predefined list of $M_i$ candidate parent sets. Each candidate set $j \in \{0, 1, \dots, M_i-1\}$ is associated with a numerical score $s_{i,j}$ and a size $d_{i,j}$ representing the number of directed edges (parents) it introduces. The objective is to select exactly one candidate parent set for every node such that the total number of directed edges across all selected sets equals a prescribed integer $K$, while maximizing the sum of the scores of the selected sets. The fixed node order guarantees that all edges point from earlier to later nodes, thereby ensuring the resulting directed graph is acyclic.

\paragraph{Decision Variables.}
We introduce binary decision variables $x_{i,j} \in \{0,1\}$ for each node $i \in \{0, \dots, N-1\}$ and each candidate index $j \in \{0, \dots, M_i-1\}$. The variable $x_{i,j}$ equals $1$ if and only if the $j$-th candidate parent set for node $i$ is selected in the network structure; otherwise, $x_{i,j} = 0$.

\paragraph{QUBO Derivation.}
\textbf{Step 1 --- Objective Function.} The original maximization objective is converted to a minimization problem by negating the objective function:
\begin{align}
    \min_{x} \; f(x) = -\sum_{i=0}^{N-1} \sum_{j=0}^{M_i-1} s_{i,j} x_{i,j}.
\end{align}

\textbf{Step 2 --- Constraints.} The problem is subject to two equality constraints:
\begin{align}
    \sum_{j=0}^{M_i-1} x_{i,j} &= 1, \quad \forall i \in \{0, \dots, N-1\}, \\
    \sum_{i=0}^{N-1} \sum_{j=0}^{M_i-1} d_{i,j} x_{i,j} &= K.
\end{align}

\textbf{Step 3 --- Penalty Terms.} We convert each constraint into a quadratic penalty term using positive weights $\lambda_1$ and $\lambda_2$.

For the first constraint (exactly one parent set per node), the penalty is:
\begin{align}
    P_1(x) &= \lambda_1 \sum_{i=0}^{N-1} \left( \sum_{j=0}^{M_i-1} x_{i,j} - 1 \right)^2.
\end{align}
Expanding the square and applying the idempotent property $x_{i,j}^2 = x_{i,j}$ for binary variables yields:
\begin{align}
    \left( \sum_{j=0}^{M_i-1} x_{i,j} - 1 \right)^2 &= \sum_{j=0}^{M_i-1} x_{i,j}^2 + \sum_{\substack{j,k=0 \\ j \neq k}}^{M_i-1} x_{i,j} x_{i,k} - 2 \sum_{j=0}^{M_i-1} x_{i,j} + 1 \nonumber \\
    &= \sum_{j=0}^{M_i-1} x_{i,j} + \sum_{\substack{j,k=0 \\ j \neq k}}^{M_i-1} x_{i,j} x_{i,k} - 2 \sum_{j=0}^{M_i-1} x_{i,j} + 1 \nonumber \\
    &= 1 - \sum_{j=0}^{M_i-1} x_{i,j} + \sum_{\substack{j,k=0 \\ j \neq k}}^{M_i-1} x_{i,j} x_{i,k}.
\end{align}
Thus, the general penalty contribution is:
\begin{align}
    P_1(x) = \lambda_1 \sum_{i=0}^{N-1} \left( 1 - \sum_{j=0}^{M_i-1} x_{i,j} + \sum_{\substack{j,k=0 \\ j \neq k}}^{M_i-1} x_{i,j} x_{i,k} \right).
\end{align}

For the second constraint (total edge count), the penalty is:
\begin{align}
    P_2(x) &= \lambda_2 \left( \sum_{i=0}^{N-1} \sum_{j=0}^{M_i-1} d_{i,j} x_{i,j} - K \right)^2.
\end{align}
Expanding the square and applying $x_{i,j}^2 = x_{i,j}$ gives:
\begin{align}
    \left( \sum_{i,j} d_{i,j} x_{i,j} - K \right)^2 &= \sum_{i,j} d_{i,j}^2 x_{i,j}^2 + \sum_{\substack{(i,j) \neq (p,q)}} d_{i,j} d_{p,q} x_{i,j} x_{p,q} - 2K \sum_{i,j} d_{i,j} x_{i,j} + K^2 \nonumber \\
    &= \sum_{i,j} d_{i,j}^2 x_{i,j} + \sum_{\substack{(i,j) \neq (p,q)}} d_{i,j} d_{p,q} x_{i,j} x_{p,q} - 2K \sum_{i,j} d_{i,j} x_{i,j} + K^2 \nonumber \\
    &= K^2 + \sum_{i=0}^{N-1} \sum_{j=0}^{M_i-1} (d_{i,j}^2 - 2K d_{i,j}) x_{i,j} + \sum_{\substack{(i,j) \neq (p,q)}} d_{i,j} d_{p,q} x_{i,j} x_{p,q}.
\end{align}
Thus, the general penalty contribution is:
\begin{align}
    P_2(x) = \lambda_2 \left( K^2 + \sum_{i=0}^{N-1} \sum_{j=0}^{M_i-1} (d_{i,j}^2 - 2K d_{i,j}) x_{i,j} + \sum_{\substack{(i,j) \neq (p,q)}} d_{i,j} d_{p,q} x_{i,j} x_{p,q} \right).
\end{align}

\textbf{Step 4 --- Combined Hamiltonian.} The full QUBO Hamiltonian $H(x)$ is the sum of the negated objective and both penalty terms:
\begin{align}
    H(x) = f(x) + P_1(x) + P_2(x).
\end{align}

\textbf{Step 5 --- Q Matrix Entries and Offset.} Collecting constant terms, linear terms, and quadratic terms yields the offset $c$ and the symmetric QUBO matrix entries $Q_{(i,j),(p,q)}$:
\begin{align}
    c &= \lambda_1 N + \lambda_2 K^2, \\
    Q_{(i,j),(i,j)} &= -s_{i,j} - \lambda_1 + \lambda_2(d_{i,j}^2 - 2K d_{i,j}), \\
    Q_{(i,j),(i,k)} &= 2\lambda_1 + 2\lambda_2 d_{i,j} d_{i,k}, \quad \text{for } j \neq k, \\
    Q_{(i,j),(p,q)} &= 2\lambda_2 d_{i,j} d_{p,q}, \quad \text{for } i \neq p.
\end{align}
The optimization problem is equivalent to $\min_{x \in \{0,1\}^n} H(x)$, where $n = \sum_{i=0}^{N-1} M_i$.

\paragraph{Penalty Weight Justification.}
The penalty weights must be chosen large enough so that any violation of the constraints incurs an energy cost strictly greater than the maximum possible improvement in the objective function. Since each variable $x_{i,j}$ is binary, flipping a single variable changes the objective value by at most $\max_{i,j} |s_{i,j}|$. Therefore, selecting
\begin{align}
    \lambda_1 > \max_{i,j} |s_{i,j}| \quad \text{and} \quad \lambda_2 > \max_{i,j} |s_{i,j}|
\end{align}
ensures that any infeasible configuration has a strictly higher energy than the best feasible configuration, forcing the global minimum to satisfy all constraints.

\paragraph{Correctness Argument.}
Minimizing $H(x)$ over $\{0,1\}^n$ recovers the optimal network structure because (i) for any feasible assignment, both penalty terms evaluate to zero, so $H(x)$ equals the negated objective function; and (ii) for any infeasible assignment, at least one constraint is violated, contributing a penalty of at least $\min(\lambda_1, \lambda_2)$ to the energy. By the choice of $\lambda_1$ and $\lambda_2$, this penalty strictly exceeds the maximum possible reduction in the objective function achievable by any single variable flip. Consequently, the global minimum of $H(x)$ cannot be infeasible, and among all feasible solutions, it must correspond to the assignment that minimizes the negated objective, i.e., maximizes the total score.

\paragraph{Illustrative Example.}
Consider a concrete instance with $N=2$ nodes. Node $0$ has $M_0=1$ candidate set, and node $1$ has $M_1=2$ candidate sets, yielding $n=3$ binary variables $x_{0,0}, x_{1,0}, x_{1,1}$. The parameters are:
\begin{align}
    s_{0,0} = 5, \; d_{0,0} = 1; \quad s_{1,0} = 3, \; d_{1,0} = 1; \quad s_{1,1} = 4, \; d_{1,1} = 2; \quad K = 2.
\end{align}
We choose $\lambda_1 = 6$ and $\lambda_2 = 6$, satisfying the penalty justification condition. The offset is $c = 6(2) + 6(2^2) = 12 + 24 = 36$. The Q matrix entries are computed as:
\begin{align}
    Q_{(0,0),(0,0)} &= -5 - 6 + 6(1^2 - 2(2)(1)) = -11 + 6(-3) = -29, \\
    Q_{(1,0),(1,0)} &= -3 - 6 + 6(1^2 - 2(2)(1)) = -9 + 6(-3) = -27, \\
    Q_{(1,1),(1,1)} &= -4 - 6 + 6(2^2 - 2(2)(2)) = -10 + 6(4 - 8) = -34, \\
    Q_{(1,0),(1,1)} &= 2(6) + 2(6)(1)(2) = 12 + 24 = 36, \\
    Q_{(0,0),(1,0)} &= 2(6)(1)(1) = 12, \quad Q_{(0,0),(1,1)} = 2(6)(1)(2) = 24.
\end{align}
The resulting Hamiltonian is:
\begin{align}
    H(x) = 36 - 29x_{0,0} - 27x_{1,0} - 34x_{1,1} + 36x_{1,0}x_{1,1} + 12x_{0,0}x_{1,0} + 24x_{0,0}x_{1,1}.
\end{align}
Evaluating $H(x)$ over all $2^3=8$ assignments, the only feasible solution satisfying the constraints is $x_{0,0}=1, x_{1,0}=1, x_{1,1}=0$ (total edges $1+1=2$). Its energy is $H(1,1,0) = 36 - 29 - 27 = -20$. All other assignments incur a penalty of at least $6$, yielding energies $\ge -14$. Thus, the global minimum is $-20$ at $x^* = (1,1,0)$, which corresponds to selecting the only candidate for node $0$ and the first candidate for node $1$, achieving the maximum feasible score of $5+3=8$.
\end{document}